\scanpage{47}
% archive foliation: 46
pour que $F_1$ ait les propriétés dites. \deleted{\illegible{2}}

$1^\circ)$ $\uline{F}_1$ doit être stable par multiplication par $\uline{B}$, i.e. \deleted{\unsure{que}} les composantes de $b_i e_\lambda$ suivant les $d_\mu$, soient $\varphi_{i\lambda\mu}\in A$, doivent être \uline{nulles}. D'où des \uline{équations}
\[
\varphi_{i\lambda\mu}=0\qquad\left[\begin{array}{l}
1\leq i\leq d\\
1\leq\lambda\leq(n-k)d\\
1\leq\mu\leq kd
\end{array}\right]
\]
\margin{Soit $G_1$ le sous schéma de $G$ défini par ces équations. On est ramené au cas ---}

$2^\circ)$ Le $B$-module quotient \deleted{$\unsure{F_1/F|M=B^k}$} $F/F_1=M$ est projectif de rang $\unsure{k|n}$. Il a une base sur $A$ de la forme $(m_i)_{1\leq i\leq kd}$,
Soient $\xi=(\xi_1-\xi_k)$ \deleted{\unsure{des}} \inserted{une suite de $k$} \unsure{éléments} de $M$. Pour qu'ils forment une base sur $B$, il f. et s. que $B^k\xrightarrow{u_\xi}M$ défini \deleted{\unsure{soit}} bijectif. Or deux $B^k$ et $M$ ont des bases, donc on peut former $\det(u_\xi)$, la condition est que $\det(u_\xi)$ soit inversible. \deleted{Donc pour que} Pour que en un pt $x\in\operatorname{Spec}(A)$, $M_x$ soit libre sur $B_x$, il faut et il suffit qu'il existe $\xi\in(M_x)^k$ tel que $\det(u_\xi)$ soit inversible dans $A_x$. On voit facilement que l'on peut aussi prendre les $\xi\in M^k$. Donc pour tt \deleted{\unsure{un}} $\xi\in M^k$, soit $U_\xi$ l'ouvert \deleted{\unsure{par}} des non-zéros de $\det(u_\xi)$. Soit $U=\bigcup U_\xi$. C'est
\margin{\unsure{préciser un th., en complétant|préciser un peu, en complétant} \unsure{tenant compte des conditions de validation encore}}
